% Heavy-Light Decomposition (HLD)

Decomposes a tree into paths where each vertex belongs to exactly one heavy path. Allows range queries and updates on paths in $O(\log^2 N)$ when paired with a Segment Tree.

\begin{lstlisting}
struct HLD {
    int n, timer;
    vector<vector<int>> adj;
    vector<int> parent, depth, heavy, head, pos;

    HLD(int n, const vector<vector<int>>& tree_adj) : 
        n(n), timer(0), adj(tree_adj), 
        parent(n + 1), depth(n + 1), heavy(n + 1, -1), head(n + 1), pos(n + 1) {}

    int dfs_sz(int u, int p, int d) {
        int size = 1, max_c_size = 0;
        parent[u] = p; depth[u] = d;
        for (int v : adj[u]) {
            if (v != p) {
                int c_size = dfs_sz(v, u, d + 1);
                size += c_size;
                if (c_size > max_c_size) {
                    max_c_size = c_size;
                    heavy[u] = v;
                }
            }
        }
        return size;
    }

    void decompose(int u, int h) {
        head[u] = h; pos[u] = ++timer;
        if (heavy[u] != -1) {
            decompose(heavy[u], h);
        }
        for (int v : adj[u]) {
            if (v != parent[u] && v != heavy[u]) {
                decompose(v, v);
            }
        }
    }

    void init(int root = 1) {
        dfs_sz(root, 0, 0);
        decompose(root, root);
    }

    template<typename SegTree>
    int query_path(int u, int v, SegTree& seg) {
        int res = 0; // Set to identity element (e.g., -INF for max, 0 for sum)
        for (; head[u] != head[v]; v = parent[head[v]]) {
            if (depth[head[u]] > depth[head[v]]) swap(u, v);
            res += seg.query(1, 1, n, pos[head[v]], pos[v]); // Modify operation accordingly
        }
        if (depth[u] > depth[v]) swap(u, v);
        res += seg.query(1, 1, n, pos[u], pos[v]); // Node values query
        // For edge values query, use: pos[heavy_child_of_u] instead of pos[u]
        return res;
    }

    template<typename SegTree>
    void update_path(int u, int v, int val, SegTree& seg) {
        for (; head[u] != head[v]; v = parent[head[v]]) {
            if (depth[head[u]] > depth[head[v]]) swap(u, v);
            seg.update(1, 1, n, pos[head[v]], pos[v], val);
        }
        if (depth[u] > depth[v]) swap(u, v);
        seg.update(1, 1, n, pos[u], pos[v], val);
    }
};
\end{lstlisting}

Initialize via \lstinline|HLD hld(n, adj)| and call \lstinline|hld.init(root)|. The \lstinline|pos| array maps vertices to a continuous segment workspace. Pass your standard \lstinline|SegTree| template into \lstinline|query_path| and \lstinline|update_path| to run path computations.
